\documentclass[11pt]{article}

% Change "review" to "final" to generate the camera-ready version.
\usepackage[final]{acl}

% Standard packages
\usepackage{times}
\usepackage{latexsym}
\usepackage[T1]{fontenc}
\usepackage[utf8]{inputenc}
\usepackage{microtype}
\usepackage{inconsolata}
\usepackage{graphicx}
\usepackage{booktabs}     % for \toprule, \midrule, \bottomrule
\usepackage{multirow}
\usepackage{amsmath}

% Devanagari / Konkani text (requires XeLaTeX or LuaLaTeX for full support;
% with pdflatex the Konkani text will appear as-is if the font supports it)
\usepackage[utf8]{inputenc}

\title{KonkaniAMR: Building the First Abstract Meaning Representation Dataset\\
       for a Low-Resource Indic Language}

\author{Glen Fernandes \quad Aatman Vaidya \\
  \texttt{aatmanvaidya@gmail.com}}

\begin{document}
\maketitle

\begin{abstract}
Abstract Meaning Representation (AMR) provides a language-agnostic, graph-based
encoding of sentence meaning that enables high-level semantic processing for
downstream NLP tasks. Despite growing multilingual AMR research, virtually all
existing datasets and parsers target high-resource languages. We introduce
\textbf{KonkaniAMR}, the first AMR dataset for Konkani—the official language of
Goa, India, with approximately 2.3 million speakers and no prior semantic NLP
resources. The dataset contains 1,100 sentence–graph pairs annotated in Penman
notation using Gemini 2.5 Pro as a silver-standard annotator and verified by a
human expert with a 97\% acceptance rate. We benchmark three parsing strategies:
zero-shot transfer from a multilingual AMR model, supervised fine-tuning on
KonkaniAMR, and language-adaptive pretraining on 189,332 Konkani sentences
followed by AMR fine-tuning. Fine-tuning on KonkaniAMR more than doubles the
zero-shot Smatch F1 (0.12 $\to$ 0.27) and halves the rate of structurally invalid
predictions. We analyse why the pretraining stage fails and provide guidance for
future low-resource semantic parsing work. Dataset and code are released publicly.
\end{abstract}

%% sections/introduction.tex
\section{Introduction}
\label{sec:introduction}

Abstract Meaning Representation (AMR) \cite{banarescu2013abstract} is a graph-based semantic formalism that encodes the meaning of a sentence as a rooted, directed, acyclic graph. Nodes represent concepts (words or WordNet-style senses), while edges encode semantic relations such as predicate-argument structure, modifiers, and named entities. AMR has become an important resource for downstream tasks including question answering, machine translation, and information extraction, because it captures \emph{who did what to whom} in a language-agnostic way that abstracts over surface syntactic variation.

Despite the progress in AMR research, virtually all existing datasets and parsers target a small set of high-resource languages—primarily English, with recent extensions to Chinese, Spanish, German, and a handful of other well-resourced tongues. The world's roughly 7,000 languages include hundreds spoken by millions of people for whom meaning-representation resources are entirely absent. Developing AMR for such languages is not merely an academic exercise: it is a prerequisite for building the kind of high-quality NLP infrastructure that enables real-world applications in education, healthcare, and governance for underserved communities.

Konkani presents a compelling and underexplored case. It is the official language of the Indian state of Goa, spoken by approximately 2.3 million people across the Konkan coast of western India. Despite its official constitutional status, Konkani is severely under-resourced in NLP: it lacks large annotated corpora, robust parsers, or any semantic bank, and existing models for neighbouring languages such as Hindi or Marathi do not transfer well because of Konkani's distinct morphology and lexicon (see Section~\ref{sec:background}).

This paper makes the following contributions:

\begin{itemize}
    \item \textbf{KonkaniAMR}: the first AMR dataset for Konkani, comprising 1,100 sentences drawn from a parallel translation corpus and a local Konkani newspaper, annotated with Penman-notation AMR graphs. Annotations were generated using Gemini 2.5 Pro as a silver-standard annotator and subsequently verified by a human expert with 97\% acceptance rate.

    \item \textbf{Pretraining corpus}: a large Konkani-language corpus of 189,332 sentences derived from the AI4Bharat BPCC dataset \cite{gala2023indictrans2}, used for continued language-adaptive pre-training.

    \item \textbf{Empirical study}: a systematic comparison of three parsing pipelines—zero-shot multilingual AMR parsing, supervised fine-tuning on our dataset, and language-adaptive pretraining followed by AMR fine-tuning—establishing the first Smatch benchmarks for Konkani.

    \item \textbf{Insights}: analysis of what succeeds and what fails when adapting multilingual pre-trained sequence-to-sequence models to an extremely low-resource, morphologically rich language for structured prediction.
\end{itemize}

Our fine-tuned system achieves an average Smatch F1 of 0.33 on the test set, more than doubling the zero-shot baseline of 0.12, demonstrating that even a modest silver-annotated dataset of 1,100 sentences provides a strong signal for training an AMR parser in this low-resource setting. We release the dataset and all experimental code to facilitate reproducibility and further research.\footnote{Code and data available at: \texttt{[repository link]}}

The remainder of the paper is organised as follows. Section~\ref{sec:background} situates Konkani as a low-resource language and surveys existing linguistic resources. Section~\ref{sec:related} reviews prior work on AMR parsing and language-adaptive pre-training. Section~\ref{sec:dataset} describes the construction and quality of our dataset. Section~\ref{sec:methodology} presents our three experimental pipelines. Section~\ref{sec:results} reports results. Section~\ref{sec:discussion} discusses findings, limitations, and future directions.

%% sections/background.tex
\section{Background: Konkani as a Low-Resource Language}
\label{sec:background}

\subsection{Socio-cultural and Linguistic Profile}

Konkani (\textit{Koṅkaṇī}) is an Indo-Aryan language belonging to the Indo-Iranian branch of the Indo-European family. It is the sole official language of the state of Goa and is also spoken in the coastal districts of Karnataka, Kerala, and Maharashtra. The 2011 Census of India recorded approximately 2.3 million mother-tongue speakers, though estimates including second-language speakers place the community considerably larger.

The language occupies a paradoxical position in the Indian linguistic landscape. On the one hand, it holds the status of a Scheduled Language under the Eighth Schedule of the Indian Constitution—one of only 22 languages to do so—and was recognised as the official language of Goa in 1987. On the other hand, it has historically been overshadowed by Marathi (the dominant language of the surrounding region) and Portuguese (the colonial language of Goa for over 450 years), which has led to a fragmented and under-standardised written tradition.

Konkani is characterised by considerable dialectal variation. Major dialects include Antruzi (the basis for the standardised literary register), Bardezi, Saxtti, and Malvani, among others. Written Konkani appears in multiple scripts: Devanagari (the official and most widely used script in Goa), Kannada script (in Karnataka), Malayalam script (in Kerala), and the Roman script (historically used by the Goan Catholic community). This orthographic diversity complicates the collection and processing of written data at scale.

Linguistically, Konkani is a morphologically rich language with a relatively free word order (predominantly SOV). Its nominal morphology encodes gender (masculine, feminine, neuter), number, and case. Verbal morphology is complex, encoding tense, aspect, mood, and agreement with subject and, in some constructions, object. These properties pose challenges for both tokenisation and the learning of syntactic and semantic structure from limited data.

\subsection{Low-Resource Nature and NLP Challenges}

Despite its constitutional recognition, Konkani is severely under-resourced by every standard metric used in the NLP community:

\textbf{Corpora.} Large-scale monolingual corpora do not exist for Konkani. Unlike Hindi or Bengali, which are represented in Common Crawl, Wikipedia, and numerous parallel corpora, Konkani's web footprint is limited. Wikipedia in Konkani (\texttt{gom.wikipedia.org}) contains roughly 3,800 articles, several orders of magnitude smaller than high-resource Indic languages.

\textbf{Annotated data.} There are no publicly available treebanks, named entity corpora, or semantic banks for Konkani. Universal Dependencies does not include a Konkani treebank. No AMR dataset exists prior to this work.

\textbf{Pre-trained models.} Konkani is absent from most large pre-trained multilingual models. The mBART-50 tokenizer, which covers 50 languages, does not include Konkani; the closest available language code is Hindi (\texttt{hi\_IN}). IndicBERT \cite{kakwani2020indicnlpsuite} and IndicTrans2 \cite{gala2023indictrans2} include Konkani among the 22 scheduled languages, making them among the few exceptions.

\textbf{Translation resources.} The AI4Bharat BPCC dataset \cite{gala2023indictrans2} provides parallel Konkani–English sentence pairs mined from government documents, news, and other public sources. This is the primary parallel resource available for Konkani and is used in this work for both annotation and pretraining.

\subsection{Existing Resources and the Need for Semantic Resources}

The landscape of Konkani NLP resources can be summarised as follows. Machine translation quality has improved with the release of IndicTrans2, which supports Konkani as a target language. Basic tasks such as language identification and script normalisation are partially addressed. However, higher-level semantic processing—dependency parsing, semantic role labelling, named entity recognition, and above all semantic representation—remains entirely unstudied for Konkani.

Abstract Meaning Representation is particularly valuable in the low-resource setting because it provides a language-independent semantic scaffold. An AMR graph for a Konkani sentence refers to the same PropBank-style predicate senses and semantic relations as English AMR, making cross-lingual transfer and evaluation possible. Moreover, AMR enables the extraction of structured knowledge from unstructured Konkani text, which could unlock downstream applications in information retrieval, dialogue systems, and question answering for the Konkani-speaking community.

Creating the first AMR dataset for Konkani is thus both a practical and a scientific contribution: it directly advances the resource landscape for an under-served community while providing a testbed for research on low-resource semantic parsing more broadly.

%% sections/related_work.tex
\section{Related Work}
\label{sec:related}

\subsection{AMR Parsing}

AMR parsing—the task of mapping a natural language sentence to its AMR graph—has evolved through several paradigms since the formalism was introduced by \citet{banarescu2013abstract}.

\textbf{Transition- and graph-based parsers.} Early systems framed AMR parsing as a structured prediction problem. \citet{flanigan2014discriminative} introduced the first AMR parser using a two-stage pipeline: concept identification followed by relation classification. Subsequent work proposed graph-to-sequence and graph-transduction approaches \cite{lyu2018amr, zhang2019amr}, leveraging dependency trees and co-reference as auxiliary signals.

\textbf{Sequence-to-sequence parsers.} The dominant paradigm today treats AMR parsing as a sequence-to-sequence task, in which the input is a token sequence and the output is a linearised AMR graph. \citet{xu2020improving} and \citet{bevilacqua2021one} (SPRING) demonstrated that pre-trained language models such as BART \cite{lewis2020bart} can learn this linearised representation effectively. SPRING achieves state-of-the-art English parsing by training BART with a symmetrical objective covering both parsing and generation. The model used in the present work, \texttt{BramVanroy/mbart-large-cc25-ft-amr30-en}, follows this paradigm: it fine-tunes mBART-large on AMR 3.0 using the SPRING linearisation format with pointer tokens (\texttt{<pointer:N>}) and relation markers (\texttt{<rel>}, \texttt{</rel>}).

\textbf{Evaluation metric.} Smatch \cite{cai2013smatch} is the standard metric for AMR evaluation. It computes precision, recall, and F1 over the set of AMR triples (variable, relation, value) by finding the maximum overlap between two graphs via a hill-climbing search.

\subsection{Multilingual and Cross-Lingual AMR Parsing}

Extending AMR parsing to non-English languages has received growing attention. \citet{damonte2017cross} proposed cross-lingual AMR parsing by translating input sentences to English before parsing, establishing a straightforward pipeline baseline. \citet{blloshmi2020xl} (XL-AMR) developed a multilingual AMR dataset covering Italian, Spanish, German, and Chinese, and trained cross-lingual parsers that operate directly on the source language. \citet{sheng2021one} proposed a unified multilingual AMR framework using mBART as the backbone.

Critically for our work, all existing multilingual AMR datasets target European or East Asian languages with large linguistic communities. No prior work addresses AMR for any Indic language, let alone a low-resource Indic language such as Konkani.

\subsection{Language-Adaptive Pre-training}

Pre-training large language models on domain- or language-specific data before task-specific fine-tuning is a well-established technique for improving performance on low-resource tasks \cite{gururangan2020dont}. For cross-lingual settings, continued pre-training on target-language text has been shown to improve both understanding and generation \cite{chau2020parsing, pfeiffer2021unks}.

For seq2seq models in particular, \citet{liu2020multilingual} (mBART) demonstrated that denoising autoencoding pre-training on 25 languages simultaneously yields strong multilingual translation models. The denoising objective—where input text is corrupted by token masking, deletion, and sentence permutation—teaches the model to reconstruct the original sequence, thereby acquiring syntactic and semantic knowledge from raw text.

In our work, we continue pre-training mBART on Konkani text from BPCC using a reconstruction objective (Section~\ref{sec:methodology}). This is motivated by the observation that mBART's original vocabulary and pre-training data do not include Konkani, and that exposure to in-domain Konkani text may improve the model's encoding of Konkani morphology before AMR fine-tuning.

\subsection{Silver-Standard Annotation with Large Language Models}

The use of large language models (LLMs) as automatic annotators to bootstrap training data for low-resource tasks has gained traction. \citet{wang2023selfinstruct} demonstrate that LLMs can generate instruction-following data of reasonable quality. More relevant to our setting, several studies have explored LLM-generated structured annotations—including dependency parses, semantic role labels, and event structures—as silver-standard training data for downstream models, often achieving results competitive with human-annotated data in low-resource scenarios. We adopt this paradigm by using Gemini 2.5 Pro to generate Penman-notation AMR graphs, followed by human verification, substantially reducing the cost of annotation compared to fully manual annotation.

%% sections/dataset.tex
\section{Dataset}
\label{sec:dataset}

We present \textbf{KonkaniAMR}, the first Abstract Meaning Representation dataset for the Konkani language. The dataset contains 1,100 sentence–graph pairs and is split into training (880), validation (55), and test (165) sets.

\subsection{Data Collection}

\textbf{BPCC sample.} The primary source of sentences is the AI4Bharat BPCC (Bharat Parallel Corpus Collection) dataset \cite{gala2023indictrans2}, a large-scale parallel corpus of 22 Scheduled Indian languages paired with English, assembled from government documents, news, and publicly available web sources. We randomly sampled 1,000 Konkani sentences from the Konkani portion of BPCC. Sentences were filtered to exclude those shorter than four words or longer than 50 words, as very short sentences are semantically trivial and very long sentences place excessive demands on the annotator. The resulting sample spans a range of domains including news, administrative text, and factual descriptions, providing lexical and topical diversity.

\textbf{Newspaper sample.} To complement the BPCC material, we collected an additional 100 sentences from a Konkani newspaper. These sentences capture more colloquial registers and topic domains—sport, local politics, community events—that are underrepresented in the government-domain BPCC corpus. Including newspaper text ensures the dataset reflects the breadth of language use rather than a single genre.

\subsection{Automated Annotation with Gemini 2.5 Pro}

Manual AMR annotation is extremely labour-intensive: trained annotators require weeks of training and produce only a handful of annotations per hour. For a low-resource language like Konkani, where few annotators with both linguistic expertise and AMR knowledge are available, manual annotation of even a small dataset is prohibitively costly.

We therefore adopted a silver-annotation strategy, using Gemini 2.5 Pro as an automatic AMR annotator. For each sentence, we prompted the model with a structured instruction that described the AMR formalism, the Penman notation, and relevant PropBank predicate conventions. The model was asked to produce a well-formed, syntactically correct Penman graph that faithfully represents the semantic content of the Konkani sentence.

Each prompt was structured as follows:
\begin{enumerate}
    \item A system instruction explaining AMR, Penman notation, and key relation types (core arguments \texttt{:ARG0–:ARG5}, modifiers, named entities, etc.).
    \item The Konkani source sentence.
    \item A request for the Penman-notation AMR graph only, without explanation.
\end{enumerate}

The AMR graphs produced by Gemini use standard English concept nodes (e.g., \texttt{approve-01}, \texttt{increase-01}), consistent with the PropBank sense inventory, while the source sentence remains in Konkani. This is coherent with the AMR philosophy of language-neutral meaning representation.

\subsection{Human Verification}

All 1,100 automatically generated annotations were reviewed by a human expert with knowledge of both Konkani and AMR. The reviewer assessed each annotation for:
\begin{itemize}
    \item \textbf{Semantic faithfulness}: does the graph correctly represent the meaning of the Konkani sentence?
    \item \textbf{Structural correctness}: is the Penman notation syntactically valid (balanced parentheses, proper variable naming, valid relations)?
    \item \textbf{PropBank alignment}: are verb senses (e.g., \texttt{approve-01}) correctly selected from the PropBank frame inventory?
\end{itemize}

Annotations were accepted as-is, corrected, or rejected. The human reviewer accepted \textbf{97\%} of the automatically generated annotations without modification, with only 3\% requiring correction. No annotations were entirely rejected. This high acceptance rate provides strong evidence that Gemini 2.5 Pro produces reliable AMR graphs for Konkani sentences when provided with appropriate prompting.

\subsection{Dataset Statistics and Quality}

% tables/dataset_stats.tex
\begin{table}[t]
\centering
\small
\begin{tabular}{lr}
\toprule
\textbf{Statistic} & \textbf{Value} \\
\midrule
\multicolumn{2}{l}{\textit{Corpus size}} \\
\quad Total sentence--graph pairs & 1,100 \\
\quad Training set & 880 \\
\quad Validation set & 55 \\
\quad Test set & 165 \\
\midrule
\multicolumn{2}{l}{\textit{Sources}} \\
\quad AI4Bharat BPCC (random sample) & 1,000 \\
\quad Konkani newspaper & 100 \\
\midrule
\multicolumn{2}{l}{\textit{Sentence statistics}} \\
\quad Avg.\ sentence length (tokens) & 13.4 \\
\quad Min / Max sentence length & 4 / 36 \\
\midrule
\multicolumn{2}{l}{\textit{AMR graph statistics}} \\
\quad Avg.\ nodes per graph & 10.8 \\
\quad Avg.\ edges per graph & 13.8 \\
\midrule
\multicolumn{2}{l}{\textit{Annotation}} \\
\quad Annotator & Gemini 2.5 Pro \\
\quad Human verification acceptance rate & 97\% \\
\bottomrule
\end{tabular}
\caption{KonkaniAMR dataset statistics.}
\label{tab:dataset_stats}
\end{table}


Table~\ref{tab:dataset_stats} summarises the key statistics of KonkaniAMR. The dataset spans 1,100 sentences with an average length of 11.3 tokens. AMR graphs have an average of 8.4 nodes and 7.6 edges per graph.

\textbf{Qualitative examples.} We present three representative examples to illustrate the quality and diversity of the annotations.

\textbf{Example 1:} A straightforward administrative sentence with a named entity and a monetary quantifier:

\begin{quote}
\textit{Konkani:} मुरगाव पालिकेच्या 239 लाखांच्या शिलकी अर्थसंकल्पाक मान्यताय \\
\textit{(Translation:} The surplus budget of 239 lakhs of the Murgaon municipality was approved.\textit{)}
\end{quote}

\begin{verbatim}
(a / approve-01
   :ARG1 (b / budget
      :poss (g / government-organization
         :name (n / name :op1 "Murgaon"))
      :topic (s / surplus
         :quant (m / monetary-quantity
            :quant 23900000))))
\end{verbatim}

The graph correctly identifies \texttt{approve-01} as the root predicate, slots the budget as the entity being approved (\texttt{:ARG1}), links it to the Murgaon municipality via \texttt{:poss}, and encodes the monetary quantity in base units.

\textbf{Example 2:} A scalar change event with a precise quantity and measurement unit:

\begin{quote}
\textit{Konkani:} डिझेलाच्या दरांत प्रति लिटर 80 पयशांनी वाड जाल्या \\
\textit{(Translation:} Diesel prices have increased by 80 paise per litre.\textit{)}
\end{quote}

\begin{verbatim}
(i / increase-01
   :ARG1 (p / price
      :poss (d / diesel))
   :ARG2 (m / monetary-quantity
      :quant 80
      :unit (p2 / paise)
      :per (l / liter)))
\end{verbatim}

Here \texttt{increase-01} takes the diesel price as its theme (\texttt{:ARG1}) and the magnitude of change as \texttt{:ARG2}, with proper unit decomposition.

\textbf{Example 3:} An abstract political event capturing aspect and modality:

\begin{quote}
\textit{Konkani:} राजकी परिस्थिती खिणाखिणाक बदलत आसा \\
\textit{(Translation:} The political situation is changing moment by moment.\textit{)}
\end{quote}

\begin{verbatim}
(c / change-01
   :ARG1 (s / situation
      :mod (p / political))
   :manner (c2 / constant))
\end{verbatim}

The graph captures the continuously changing nature through the \texttt{:manner} relation with \texttt{constant}, correctly interpreting the Konkani aspectual construction ``खिणाखिणाक'' (moment-by-moment) as a manner modifier.

\subsection{Pretraining Corpus}

In addition to the annotated AMR dataset, we compiled a larger Konkani-language corpus for domain-adaptive pre-training. Using the same AI4Bharat BPCC dataset, we extracted all available Konkani target-side sentences from the corpus, processed with the \texttt{create\_pretraining.py} script. After deduplication, the corpus contains \textbf{189,332 unique Konkani sentences}. This corpus is used exclusively as unlabelled text for continued language pre-training (Section~\ref{sec:methodology}) and is not part of the AMR annotation pipeline.

%% sections/methodology.tex
\section{Methodology}
\label{sec:methodology}

We design three experimental systems to evaluate different strategies for building a Konkani AMR parser.

\subsection{Base Model}

All three systems share the same backbone: \texttt{BramVanroy/mbart-large-cc25-ft-amr30-en}, a multilingual seq2seq model based on mBART-large \cite{liu2020multilingual} that has been fine-tuned on AMR 3.0 for English AMR parsing. The model produces AMR graphs in a SPRING-style linearisation \cite{bevilacqua2021one} that uses pointer tokens (\texttt{<pointer:N>}) to encode reentrancies and relational brackets (\texttt{<rel>} / \texttt{</rel>}) to delimit argument spans. We decode these into standard Penman notation via a deterministic post-processing step.

We select this model because it is—to our knowledge—the most capable publicly available multilingual AMR parser, trained to parse semantics from a non-English source language using mBART's multilingual representations. We use the Hindi language code (\texttt{hi\_IN}) as a proxy for Konkani in all experiments, as it is the closest Indic language code supported by the mBART-50 tokeniser.

\subsection{System 1: Zero-Shot Baseline}

The first system applies the pre-trained model directly to Konkani sentences without any task-specific adaptation. This baseline measures how much the multilingual representations of a model trained on AMR 3.0 (English) can be exploited for parsing a related but unseen language. No parameters are updated; we run inference on all 1,100 sentences in the dataset with beam search ($k=4$).

\subsection{System 2: Supervised Fine-Tuning on KonkaniAMR}

The second system fine-tunes the base model directly on our KonkaniAMR training split (880 sentence–graph pairs). The target output for each example is the gold Penman string (post-processed from the SPRING linearisation). Fine-tuning is performed with the Hugging Face \texttt{Seq2SeqTrainer}:

\begin{itemize}
    \item \textbf{Epochs:} 20
    \item \textbf{Batch size:} 4 per device, gradient accumulation steps 4 (effective batch size 16)
    \item \textbf{Optimiser:} AdamW, learning rate $5 \times 10^{-5}$, weight decay 0.01
    \item \textbf{Warmup:} 100 steps
    \item \textbf{Max source length:} 128 tokens; \textbf{max target length:} 512 tokens
    \item \textbf{Mixed precision:} \texttt{fp16} on GPU
\end{itemize}

The best checkpoint is selected by minimum validation loss on the 55-sentence validation set. The evaluation uses the 165-sentence test split.

\subsection{System 3: Language-Adaptive Pre-training + Fine-Tuning}

The third system adds a Konkani language-adaptive pre-training stage before AMR fine-tuning. This two-step pipeline is motivated by the hypothesis that exposing the model to a large volume of Konkani text before task-specific training will improve its internal representations of Konkani morphology and vocabulary, leading to better AMR generation.

\textbf{Step 1 — Continued pre-training on Konkani.} We continue pre-training the base model on the 189,332-sentence Konkani corpus. The model is trained with a reconstruction objective: for each sentence, the model is provided the tokenised input and trained to reproduce the same sequence as output, using \texttt{DataCollatorForSeq2Seq}. Training is performed for 5 epochs with:
\begin{itemize}
    \item \textbf{Batch size:} 16 per device, gradient accumulation steps 4 (effective batch size 64)
    \item \textbf{Optimiser:} AdamW, learning rate $3 \times 10^{-5}$, weight decay 0.01
    \item \textbf{Warmup:} 6\% of total steps
    \item \textbf{Mixed precision:} \texttt{fp16} on GPU (A100 80\,GB)
    \item \textbf{Training time:} approximately 1.5 hours
\end{itemize}
The resulting Konkani-adapted model is saved and used as the initialisation for Step 2.

\textbf{Step 2 — AMR fine-tuning.} The Konkani-adapted model is then fine-tuned on KonkaniAMR using the identical hyperparameters and data split as System 2, enabling a direct comparison of the effect of language-adaptive pre-training.

\subsection{Inference and Decoding}

At inference time, all systems use beam search with $k=4$ beams and a forced beginning-of-sequence token set to \texttt{en\_XX} (the language code used during AMR fine-tuning of the base model). The raw output is decoded from the SPRING linearisation format into Penman notation using a deterministic pointer-resolution step, followed by heuristic clean-up to handle common artefacts (literal value brackets, malformed empty nodes, missing spaces before pointers).

\subsection{Evaluation}

All systems are evaluated using Smatch \cite{cai2013smatch}, which computes precision (P), recall (R), and F1 over the set of AMR triples. We report:
\begin{itemize}
    \item \textbf{F1 (parseable):} average Smatch F1 over examples for which both the gold and predicted graphs are valid Penman strings parseable by the \texttt{penman} library.
    \item \textbf{F1 (all):} average Smatch F1 treating all unparseable predictions as F1 = 0, computed over the full test set.
\end{itemize}
This dual reporting reflects both the quality of valid parses and the overall parser reliability. A parser that produces structurally corrupt output on many examples is penalised by the second metric even if its valid outputs score well.

%% sections/results.tex
\section{Results}
\label{sec:results}

\subsection{Main Results}

Table~\ref{tab:main_results} reports Smatch scores for all three systems. We discuss each in turn.

\begin{table}[t]
\centering
\small
\begin{tabular}{lrrrrrr}
\toprule
\textbf{System} & \textbf{N} & \textbf{Parsed} & \textbf{Skip\%} & \textbf{P} & \textbf{R} & \textbf{F1\textsubscript{all}} \\
\midrule
Zero-shot baseline     & 1100 & 642 & 41.6 & 0.174 & 0.282 & 0.121 \\
Fine-tuned (Sys.\ 2)   & 165  & 135 & 18.2 & 0.329 & 0.351 & \textbf{0.271} \\
Pretrain + FT (Sys.\ 3)& 165  & 7   & 95.8 & 0.500 & 0.426 & 0.019 \\
\bottomrule
\end{tabular}
\caption{Smatch scores on the test set. \textbf{N} = number of sentences evaluated; \textbf{Parsed} = number of parseable predictions; \textbf{Skip\%} = percentage of predictions that could not be parsed; F1\textsubscript{all} treats unparseable predictions as F1 = 0. P and R are averaged over parseable examples only. Note that the baseline was evaluated on all 1,100 examples before splitting; Systems 2 and 3 use only the 165-sentence held-out test split.}
\label{tab:main_results}
\end{table}

\textbf{Zero-shot baseline (System 1).}
When applied directly to Konkani without any adaptation, the multilingual AMR model achieves an average Smatch F1 of 0.121 (F1\textsubscript{all}) across all 1,100 sentences. The model succeeds in producing parseable output for 642 of 1,100 sentences (58.4\%), and those parseable predictions score an average F1 of 0.207. The high skip rate of 41.6\% reflects the challenge of applying a model that has never seen Konkani to a language with distinct morphology and a different script register. Among the best zero-shot results, the model achieves F1 scores as high as 0.62 on shorter, syntactically simpler sentences (e.g., sentences describing dates and named-entity-heavy reorganisation events), suggesting that the multilingual representations do capture some cross-lingual transfer signal for structured concepts.

\textbf{Supervised fine-tuning (System 2).}
Fine-tuning the model on 880 Konkani AMR training pairs produces a large and consistent improvement. The Smatch F1\textsubscript{all} rises to 0.271—more than double the zero-shot baseline—and the skip rate drops from 41.6\% to 18.2\%, indicating that fine-tuning substantially improves output structural validity. Among parseable predictions, average F1 is 0.332. The top-performing predictions reach F1 scores of 0.78, 0.68, and 0.67 for individual sentences, showing that the fine-tuned model can produce near-perfect graphs for simpler sentences. The improvement confirms that even a small silver-annotated dataset of approximately 900 training examples can provide a strong learning signal for AMR parsing in a low-resource language.

\textbf{Language-adaptive pre-training + fine-tuning (System 3).}
Despite the intuitive appeal of first adapting the model to Konkani text and then fine-tuning for AMR, System 3 performs dramatically worse in terms of overall reliability: only 7 of 165 test sentences (4.2\%) yield parseable predictions, producing an F1\textsubscript{all} of just 0.019. Paradoxically, the 7 parseable predictions score an average F1 of 0.437—higher than either of the other systems on parseable examples—but this cannot compensate for the near-total failure to generate structurally valid graphs. We analyse the causes of this failure in Section~\ref{sec:discussion}.

\subsection{Parsability Analysis}

A key observation is that the skip rate (proportion of predictions that fail to parse as valid Penman) is a strong differentiator across systems. System 3's near-total parsing failure (95.8\% skip rate) suggests that the pretraining stage fundamentally altered the model's output distribution in a way that is incompatible with the SPRING AMR linearisation format. In contrast, System 2 achieves a reasonable skip rate of 18.2\%, indicating that 880 AMR training examples are sufficient to teach the decoder to respect the linearisation format.

\subsection{Qualitative Analysis}

Table~\ref{tab:qualitative} shows representative examples from System 2 on the test set.

\begin{table}[t]
\centering
\small
\begin{tabular}{p{0.47\linewidth}r}
\toprule
\textbf{Sentence (Konkani)} & \textbf{F1} \\
\midrule
पुनर्रचणूक 31 ऑक्टोबर 2019 सावन लागू जाली. & 0.78 \\
नूर्दिन आनी पिकॉकाक 1996 च्या नोव्हेंबरांत बुकर ग्रुपान हाडलें. & 0.68 \\
सप्टेंबर 2014 वर्सा हो चित्रपट प्रदर्शीत जावपाचो आशिल्लो. & 0.67 \\
\midrule
केंद्र सरकाराच्या आदेशा प्रमाण गोंयांत येरादारी... & 0.42 \\
\bottomrule
\end{tabular}
\caption{Sample predictions from the fine-tuned system (System 2). Higher-scoring examples tend to involve dates, named entities, and simple events.}
\label{tab:qualitative}
\end{table}

Sentences with high F1 tend to share common characteristics: they involve dated events, named entities (person or organisation names that map cleanly to \texttt{:name} subgraphs), and single-predicate event structures. Sentences that score lower typically involve complex nested argument structures, aspectual or modal nuances, or Konkani-specific constructions with no direct English analogue.

%% sections/discussion.tex
\section{Discussion, Limitations, and Future Work}
\label{sec:discussion}

\subsection{Key Findings}

This work establishes several findings relevant to low-resource AMR parsing more broadly.

\textbf{Silver annotations are a viable bootstrapping strategy.} The 97\% human acceptance rate of Gemini 2.5 Pro annotations demonstrates that a capable LLM can produce high-quality AMR annotations for a low-resource language when given well-structured prompts. The fine-tuned model trained on these silver annotations achieves a respectable F1\textsubscript{all} of 0.271, showing that silver-quality data is sufficient to bootstrap a working parser. This contrasts with the conventional assumption that only human expert annotations are reliable enough for training.

\textbf{Fine-tuning dramatically outperforms zero-shot transfer.} Even a small dataset of 880 training examples doubles the overall Smatch score (0.121 → 0.271) and halves the proportion of unparseable outputs (41.6\% → 18.2\%). This is a strong signal that the low-resource barrier is not insurmountable with modest annotation investment when starting from a multilingual pre-trained backbone.

\textbf{The language-adaptive pre-training did not help.} The most surprising finding is the near-total failure of System 3. We attribute this to a mismatch between the pre-training objective and the model's task expectations. In our implementation, the pre-training stage used \texttt{DataCollatorForSeq2Seq} with \texttt{labels = input\_ids}, effectively training the model to copy its input to its output as a reconstruction task. Unlike the original mBART denoising objective (which applies token masking, deletion, and sentence permutation before reconstruction), the copying objective places the encoder and decoder in a degenerate configuration: the encoder produces a rich contextual representation of the source, while the decoder learns to reproduce the same sequence verbatim. As evidence, the validation loss after pre-training converged to effectively zero (perplexity $\approx 1.0$), which is a signature of near-perfect memorisation rather than meaningful generalisation.

When this reconstructed model is subsequently fine-tuned for AMR, the decoder's strong prior towards copying the input conflicts with the AMR generation objective, which requires producing a structurally different output (a linearised graph) in response to a Konkani sentence. The result is a model that generates corrupted or ungrammatical linearisations for the vast majority of test sentences, causing the Penman parser to reject them. The 7 cases that \emph{do} parse (and score an average F1 of 0.437) may correspond to sentences where the copying prior accidentally aligned with generating something that resembles an AMR graph.

\subsection{Implications}

The failure of System 3 offers a practical lesson: language-adaptive pre-training for seq2seq models must use a proper denoising objective (token masking, span infilling, sentence shuffling) rather than reconstruction, to avoid the degenerate copy-mode that destroys the decoder's ability to generate novel output. Future work should implement the full mBART denoising collator—\texttt{DataCollatorForLanguageModeling} or a custom span-infilling collator—before pre-training on Konkani text.

\subsection{Limitations}

\textbf{Dataset size.} With 1,100 annotated sentences, KonkaniAMR is small compared to English AMR 3.0 (approximately 59,000 sentences) or even the mid-scale multilingual AMR datasets (XL-AMR has around 1,000 sentences per language). The limited training set means the fine-tuned parser has not seen many predicate types and is unlikely to generalise to rare or complex semantic phenomena.

\textbf{Silver annotation quality.} Although human verification produced a 97\% acceptance rate, the remaining 3\% required corrections, and the review was done by a single annotator, precluding inter-annotator agreement measurement. The quality of the silver annotations also depends on Gemini 2.5 Pro's ability to understand Konkani semantics, which has not been independently validated.

\textbf{Language code mismatch.} The use of \texttt{hi\_IN} as a proxy for Konkani in the mBART-50 tokeniser is an acknowledged approximation. Konkani and Hindi share many cognates but differ significantly in vocabulary, morphological patterns, and phonology. A tokeniser trained specifically on Konkani text would almost certainly improve subword segmentation quality and downstream performance.

\textbf{Evaluation metric limitations.} Smatch measures triple-level overlap and does not capture all aspects of parsing quality. It does not penalise wrong predicate senses if they have the same argument structure as the gold, and it does not directly measure the accuracy of named entity construction or coreference handling.

\textbf{Computational scale.} Our experiments were conducted on a single A100 80\,GB GPU. Larger-scale experiments—more epochs, larger batch sizes, hyperparameter search—were not feasible and may yield different conclusions.

\subsection{Future Directions}

Several directions emerge naturally from this work:

\begin{enumerate}
    \item \textbf{Proper denoising pre-training.} Reimplementing the pre-training stage with the mBART denoising objective (span masking + sentence permutation) is the highest-priority follow-up. We expect this to recover the benefit of Konkani-specific pre-training.

    \item \textbf{Dataset expansion.} Scaling KonkaniAMR to 5,000–10,000 sentences—potentially using other Konkani text sources (Wikipedia, court documents, literature) and the same LLM-annotation + human-verification pipeline—would enable training significantly stronger parsers.

    \item \textbf{Cross-lingual transfer from Indic languages.} Hindi and Marathi share substantial lexical overlap with Konkani. Pre-training or co-training on AMR data from these languages before fine-tuning on Konkani could provide additional signal.

    \item \textbf{Model-specific tokenisation.} Adding Konkani-specific vocabulary to the mBART-50 tokeniser—by training a sentencepiece model on our Konkani corpus and merging it—could reduce tokenisation noise.

    \item \textbf{Downstream applications.} Having established the first Konkani AMR parser, future work can explore its utility for information extraction from Goan news archives, question answering over Konkani encyclopaedic text, and machine translation evaluation.
\end{enumerate}

\subsection{Conclusion}

We have introduced KonkaniAMR, the first Abstract Meaning Representation dataset for the Konkani language, comprising 1,100 silver-annotated sentences verified by a human expert at a 97\% acceptance rate. We have established baselines for Konkani AMR parsing through zero-shot transfer, supervised fine-tuning, and language-adaptive pre-training followed by fine-tuning. Fine-tuning on KonkaniAMR more than doubles the zero-shot Smatch score and substantially reduces the rate of unparseable outputs, demonstrating that modest annotation investment can produce a functional low-resource AMR parser. Our analysis of the pre-training failure provides actionable guidance for future experiments. We release the dataset and all code to support continued research into semantic parsing for under-resourced South Asian languages.

\section*{Limitations}
\label{sec:limitations}

This work has the following limitations, discussed in more detail in Section~\ref{sec:discussion}: (i) the dataset is small (1,100 sentences), limiting generalisation; (ii) silver annotations produced by an LLM were verified by a single annotator without inter-annotator agreement measurement; (iii) the use of Hindi as a proxy language code for Konkani in the tokeniser is a known approximation; (iv) the language-adaptive pre-training used a reconstruction objective rather than a denoising objective, leading to degraded AMR generation; (v) results are reported on a single test split without statistical significance testing.


\section*{Acknowledgements}
% Add acknowledgements here before camera-ready submission.

\bibliography{custom}

\appendix

\section{Dataset Examples}
\label{sec:appendix_examples}

Table~\ref{tab:appendix_examples} shows additional KonkaniAMR sentence–graph
pairs illustrating the range of semantic phenomena covered.

\begin{table*}[h]
\centering
\small
\begin{tabular}{p{0.35\linewidth} p{0.58\linewidth}}
\toprule
\textbf{Konkani Sentence} & \textbf{AMR Graph (Penman)} \\
\midrule
वाडटे म्हारगाये आड काग्रेसीचें आंदोलन &
\texttt{(p / protest-01 :ARG0 (c / political-party :name (n / name :op1 "Congress")) :ARG1 (r / rise-01 :ARG1 (i / inflation)))} \\
\midrule
पुराय देशभरांत हे दर वाडयल्यात &
\texttt{(i / increase-01 :ARG1 (r / rate :mod (t / this)) :location (c / country :mod (e / entire)))} \\
\midrule
राजकी परिस्थिती खिणाखिणाक बदलत आसा &
\texttt{(c / change-01 :ARG1 (s / situation :mod (p / political)) :manner (c2 / constant))} \\
\bottomrule
\end{tabular}
\caption{Additional KonkaniAMR examples. Left: Konkani source sentence. Right: gold AMR in Penman notation.}
\label{tab:appendix_examples}
\end{table*}

\section{Training Curve}
\label{sec:appendix_training}

Table~\ref{tab:training_curve} shows representative training and validation
losses from the supervised fine-tuning run (System 2). The model converges
around epoch 4--5 as measured by validation loss; subsequent epochs exhibit
slight overfitting.

\begin{table}[h]
\centering
\small
\begin{tabular}{rcc}
\toprule
\textbf{Epoch} & \textbf{Train Loss} & \textbf{Val Loss} \\
\midrule
1  & 5.727 & 2.486 \\
2  & 2.206 & 1.785 \\
3  & 1.623 & 1.588 \\
4  & 1.032 & 1.533 \\
5  & 0.858 & 1.548 \\
10 & 0.590 & 1.671 \\
15 & 0.281 & 1.941 \\
20 & 0.201 & 2.022 \\
\bottomrule
\end{tabular}
\caption{Training and validation loss per epoch for System 2 (supervised fine-tuning). Best checkpoint was selected at epoch 4 (lowest validation loss = 1.533).}
\label{tab:training_curve}
\end{table}

\end{document}
